(1) 有效势为 $V_{eff}(R)=\Phi(R)+h^2/(2R^2)$。圆轨道要求 $V_{eff}'=0$，即 $\Phi'(R_0)=h^2/R_0^3$。令 $\Omega=h/R_0^2$，得到 $\Omega^2=\Phi'(R_0)/R_0$，是向心加速度平衡式。

(2) 写 $R=R_0+\xi$ 并保留一阶，得 $\ddot\xi+\kappa^2\xi=0$，其中 $\kappa^2=V_{eff}''(R_0)=\Phi''(R_0)+3h^2/R_0^4$。由 $\Phi'=R\Omega^2$ 对 $R$ 求导，化为 $\kappa^2=R\,d\Omega^2/dR+4\Omega^2$。

(3) 计算 $d(R^4\Omega^2)/dR=R^3(R\,d\Omega^2/dR+4\Omega^2)=R^3\kappa^2$。因此 $\kappa^2>0$ 等价于 $d(h_c^2)/dR>0$，这里 $h_c=R^2\Omega(R)$ 是各半径圆轨道所需的比角动量；这是径向线性稳定的 Rayleigh 形式。

(4) 令 $\Omega=C R^{-q}$，则 $R\,d\Omega^2/dR=-2q\Omega^2$，故 $\kappa^2=(4-2q)\Omega^2$。对于非退化圆轨道，$q<2$ 稳定、$q>2$ 不稳定、$q=2$ 边界中性；开普勒引力 $q=3/2$ 给 $\kappa=\Omega$。
